
\subsubsection{6.1 Zero-Entropy Entity-Errors as Precision Errors}

Sixty percent of entity-error samples exhibited exactly zero attention entropy over the correct entity span. This pattern indicates that the model deterministically attends elsewhere rather than diffusing attention across the sequence. Entity-substitution failures are precision errors (focused-but-wrong attention) rather than recall errors (absence of attention).

This interpretation has implications for correction strategy. RAG retrieval targets entity-level misdirection by injecting the correct entity into context, enabling the model to redirect attention. COT prompting, which decomposes reasoning steps, does not address entity-level attention errors and is therefore hypothesized to be less effective for entity-substitution failures. However, this hypothesis is not empirically validated; only synthetic structural validation was performed.

\subsubsection{6.2 Model-Specific Validation}

Attention pattern detection and classification were validated on GPT-2 only (124M parameters, 12 layers). Attention patterns may differ in larger models (e.g., Llama-2-7B with 7B parameters and 32 layers, or GPT-3.5). The optimal threshold (0.32) may be model-specific.

Multi-model replication is required to determine whether the attention entropy signature generalizes across architectures and model sizes. This limitation restricts claims to GPT-2 only.

\subsubsection{6.3 Tokenizer Mismatch and Sample Loss}

Twenty-seven percent of non-entity samples were excluded due to unmappable entity spans between spaCy (character-level) and GPT-2 BPE (subword-level). Character-to-token mapping returns \texttt{None} when entity spans fragment across subword tokens.

Conservative exclusion preserves validity (only aligned spans analyzed) but reduces statistical power. The pattern remained robust (p < 0.001) despite sample loss. Future work could implement sub-word-aware span alignment to recover lost samples.

\subsubsection{6.4 Correction Effectiveness: Unvalidated}

The correction routing component (h-m2) used mock implementations with configurable synthetic success rates. No real Wikipedia API retrieval was performed. No GPT-3.5 or Llama-2 generation was performed. No GPT-judge evaluation was performed.

The experiment demonstrates that the pipeline structure is viable (data loading, dual-model framework, matched versus mismatched comparison, gate checking all function correctly). However, correction effectiveness results (RAG 52\% versus COT 28\% for GPT-3.5) are synthetic placeholders.

Real-world validation requires:

\begin{enumerate}
\item Implementing Wikipedia API retrieval (top-3 articles per entity).
\item Implementing GPT-3.5 generation with retrieved context.
\item Implementing GPT-judge for semantic equivalence evaluation.
\item Measuring actual correction success rates on 50 entity-errors.
\end{enumerate}

The claim that matched routing improves correction effectiveness by $\geq 20$ percentage points is not empirically supported and remains a hypothesis pending future validation.

\subsubsection{6.5 Manual Labeling Bottleneck}

The framework requires manual gold labels (entity-error versus non-entity-error) for the initial 100-sample dataset. Annotation requires 2-4 hours of human effort. Automated labeling is not validated.

Scalability to thousands of failures requires supervised classifier training on the 100 gold-labeled samples, with features including entropy, question length, entity count, and question type. Automated labeling would enable deployment without per-instance human annotation. This remains future work.

\subsubsection{6.6 Single Failure Type}

Pattern detection and routing were validated for entity-substitution errors only. Reasoning errors, knowledge gaps, and hybrid failure modes were not tested. The framework does not address multi-class classification or multi-failure-type routing.

Extension to multiple failure types requires:

\begin{enumerate}
\item Expanded gold-label taxonomy (entity-substitution, reasoning-error, knowledge-gap).
\item Multi-class entropy-based classifier.
\item Matched routing for each failure type.
\end{enumerate}

Generalization beyond entity-substitution is not demonstrated and remains future work.
